\documentclass[aspectratio=169, 10pt]{beamer}
\usepackage{sunass_estilo}

\title{Módulo IV: Equilibrio General y Bienestar}
\subtitle{Guía Teórica de Repaso Rápido · PEU SUNASS}
\author[PEU SUNASS · Economía]{Preparación Examen de Admisión PEU SUNASS}
\institute{Basado en el Temario Oficial de Economía · Microeconomía}
\date[2026]{2026}

\begin{document}

\begin{frame}[plain]
  \titlepage
\end{frame}

% ==============================================================================
% SLIDE 2: MAPA CONCEPTUAL
% ==============================================================================
\begin{frame}{Mapa Conceptual: Equilibrio General y Bienestar}
  \vspace{0.15cm}
  \begin{columns}[T]
    \begin{column}{0.32\textwidth}
      \tarjeta{azulPrincipal}{1. Equilibrio General}{
        Analiza todos los mercados simultáneamente (interdependencia).\\
        • \textbf{Ley de Walras:} \(\sum_{i=1}^k p_i z_i(p) \equiv 0\).\\
        • Solo importan los \textbf{precios relativos}.
      }
    \end{column}
    \begin{column}{0.32\textwidth}
      \tarjeta{naranjaAcento}{2. Caja de Edgeworth}{
        Economía de intercambio puro (\(2\times 2\)):\\
        • Eficiencia en el intercambio (\(TMS^A = TMS^B\)).\\
        • \textbf{Curva de Contrato} y Núcleo.
      }
    \end{column}
    \begin{column}{0.32\textwidth}
      \tarjeta{verdeAcento}{3. Teoremas del Bienestar}{
        • \textbf{1er Teorema:} Todo eq. competitivo es óptimo de Pareto.\\
        • \textbf{2do Teorema:} Todo óptimo de Pareto es alcanzable vía transferencias de suma alzada.
      }
    \end{column}
  \end{columns}

  \vspace{0.2cm}
  \recuadroclave{
    \small\textbf{Ley de Walras:} Si de \(k\) mercados hay \(k-1\) en equilibrio (\(z_i=0\)), el \(k\)-ésimo mercado también está necesariamente en equilibrio.
  }
  \mensajefinal{En equilibrio general solo se determinan precios relativos (\(k-1\) ecuaciones independientes para \(k\) precios); siempre se fija un numerario.}
\end{frame}

% ==============================================================================
% SLIDE 3: CAJA DE EDGEWORTH DE INTERCAMBIO (MEJORADA Y COMPLETA)
% ==============================================================================
\begin{frame}{1. La Caja de Edgeworth: Lente, Núcleo y Equilibrio}
  \begin{columns}[T]
    \begin{column}{0.54\textwidth}
      {\small\textbf{Lectura Dual (\(O_A\) abajo-izq. / \(O_B\) arriba-der.)}}
      \vspace{0.18cm}

      \begin{tikzpicture}[scale=0.80]
        % Relleno exacto del lente de mejora paretiana entre U_A^W y U_B^W (x in [1.45, 5.75])
        % U_A = x^0.6 y^0.4 -> yA(x) = 6.0238 / x^1.5
        % U_B = (6.4-x)^0.4 (4.5-y)^0.6 -> yB(x) = 4.5 - 3.0497 / (6.4-x)^(2/3)
        \fill[verdeAcento!16]
          plot[domain=1.45:5.75, samples=45] (\x, {6.0238/(\x^(1.5))}) --
          plot[domain=5.75:1.45, samples=45] (\x, {4.5 - 3.0497/((6.4-\x)^(0.6667))}) -- cycle;

        % Marco de la Caja (W=6.4, H=4.5)
        \draw[thick, aceroAcento] (0,0) rectangle (6.4,4.5);

        % Orígenes y 4 ejes con flechas direccionales claras
        \fill[azulPrincipal] (0,0) circle (2.2pt) node[anchor=north east, font=\scriptsize\bfseries] {$O_A$};
        \fill[rojoAcento] (6.4,4.5) circle (2.2pt) node[anchor=south west, font=\scriptsize\bfseries] {$O_B$};
        \node[font=\tiny\bfseries, azulPrincipal] at (3.2, -0.30) {Bien $x_1^A \;\longrightarrow$};
        \node[font=\tiny\bfseries, azulPrincipal, rotate=90] at (-0.34, 2.25) {Bien $x_2^A \;\longrightarrow$};
        \node[font=\tiny\bfseries, rojoAcento] at (3.2, 4.78) {$\longleftarrow\; \text{Bien } x_1^B$};
        \node[font=\tiny\bfseries, rojoAcento, rotate=90] at (6.72, 2.25) {$\longleftarrow\; \text{Bien } x_2^B$};

        % Curva de contrato completa y = 18*x / (57.6 - 5*x) y Núcleo destacado [2.89, 4.61]
        \draw[verdeAcento!75, thick, dashed, domain=0:6.4, samples=45] plot (\x, {18*\x/(57.6 - 5*\x)});
        \draw[verdeAcento, line width=2.2pt, domain=2.89:4.61, samples=25] plot (\x, {18*\x/(57.6 - 5*\x)});

        % Curvas de indiferencia que pasan por W(1.45, 3.45)
        \draw[azulPrincipal, thick, domain=1.22:5.95, samples=45] plot (\x, {6.0238/(\x^(1.5))});
        \node[azulPrincipal, font=\tiny\bfseries, anchor=west] at (1.35, 4.18) {$U_A^W$};
        \draw[rojoAcento, thick, domain=0.95:5.85, samples=45] plot (\x, {4.5 - 3.0497/((6.4-\x)^(0.6667))});
        \node[rojoAcento, font=\tiny\bfseries, anchor=north] at (0.92, 3.44) {$U_B^W$};

        % Curvas tangentes en E*(3.821, 1.787)
        \draw[azulPrincipal!85, thick, dashed, domain=2.15:5.95, samples=40] plot (\x, {13.3469/(\x^(1.5))});
        \node[azulPrincipal, font=\tiny, anchor=south] at (5.75, 0.92) {$U_A^*$};
        \draw[rojoAcento!85, thick, dashed, domain=0.95:5.18, samples=40] plot (\x, {4.5 - 5.1022/((6.4-\x)^(0.6667))});
        \node[rojoAcento, font=\tiny, anchor=north] at (1.08, 2.75) {$U_B^*$};

        % Recta presupuestaria por W(1.45,3.45) y E*(3.821,1.787), pendiente -0.7014
        \draw[grisTexto!90!black, thick] (0.35, 4.221) -- (5.95, 0.294);

        % Dotación W y Equilibrio E*
        \fill[naranjaAcento] (1.45, 3.45) circle (2.6pt) node[anchor=south west, font=\scriptsize\bfseries] {$W$};
        \fill[azulOscuro] (3.821, 1.787) circle (2.7pt);
        \node[azulOscuro, font=\scriptsize\bfseries, anchor=east] at (3.62, 1.56) {$E^*$};

        % Anotaciones limpias dentro/al borde sin chocar
        \node[verdeAcento!85!black, font=\tiny\bfseries, align=center] at (2.25, 2.02) {Lente\\(mejora)};
        \draw[<-, verdeAcento, thin] (4.42, 2.38) -- (5.05, 2.65) node[right, font=\tiny\bfseries, align=left] {Núcleo\\(en CC)};
        \draw[<-, verdeAcento!80!black, thin] (5.55, 3.45) -- (4.70, 3.92) node[left, font=\tiny] {Curva de Contrato};
      \end{tikzpicture}
    \end{column}

    \begin{column}{0.43\textwidth}
      \indicador{Condición de Eficiencia (\(\text{CC}\))}{verdeAcento}{\(TMS_{1,2}^A = TMS_{1,2}^B\)}
      {\scriptsize
        • En \(W\), \(TMS^A \neq TMS^B\): se forma el \textbf{lente} entre \(U_A^W\) y \(U_B^W\) con ganancias mutuas.\\
        • \textbf{Curva de Contrato:} todos los Óptimos de Pareto (\(O_A \to O_B\)).\\
        • \textbf{Núcleo:} tramo de la CC dentro del lente (ambos ganan vs. \(W\)).\\[3pt]
      }
      \indicador{Equilibrio Walrasiano (\(E^*\))}{azulPrincipal}{\(TMS^A = \frac{p_1}{p_2} = TMS^B\)}
      {\scriptsize
        • La recta de precios (pendiente \(-\frac{p_1}{p_2}\)) parte de \(W\) y es tangente en \(E^* \in \text{Núcleo}\).\\
        • En \(E^*\), \(z_1 = z_2 = 0\) (excesos de demanda nulos).
      }
    \end{column}
  \end{columns}

  \mensajefinal{Dentro del lente ambos ganan frente a \(W\); solo la tangencia sobre la CC es Óptimo de Pareto.}
\end{frame}

% ==============================================================================
% SLIDE 4: CASO NUMÉRICO PEU SUNASS Y LEY DE WALRAS EN LA CAJA
% ==============================================================================
\begin{frame}{1b. Caso Numérico PEU SUNASS y Desequilibrio Walrasiano}
  \begin{columns}[T]
    \begin{column}{0.51\textwidth}
      {\small\textbf{Economía \(U_A = x_A y_A,\; U_B = x_B y_B\), \(W: A=(8,2)\)}}
      \vspace{0.12cm}

      \begin{tikzpicture}[xscale=0.54, yscale=0.45]
        % Lente exacto entre W=(8,2) y (2,8): U_A=16 -> 16/x; U_B=16 -> 10 - 16/(10-x)
        \fill[verdeAcento!15]
          plot[domain=2:8, samples=40] (\x, {16/\x}) --
          plot[domain=8:2, samples=40] (\x, {10 - 16/(10-\x)}) -- cycle;

        % Caja 10x10
        \draw[thick, aceroAcento] (0,0) rectangle (10,10);
        \fill[azulPrincipal] (0,0) circle (3pt) node[anchor=north east, font=\scriptsize\bfseries] {$O_A$};
        \fill[rojoAcento] (10,10) circle (3pt) node[anchor=south west, font=\scriptsize\bfseries] {$O_B$};

        % Marcas en ejes A
        \draw[grisTexto!50, dotted] (8,0) -- (8,2) -- (0,2);
        \draw[grisTexto!50, dotted] (5,0) -- (5,5) -- (0,5);
        \node[below, font=\tiny] at (5,0) {$5$};
        \node[below, font=\tiny] at (8,0) {$8$};
        \node[below right, font=\tiny, grisTexto] at (9.6,0) {$X=10$};
        \node[left, font=\tiny] at (0,2) {$2$};
        \node[left, font=\tiny] at (0,5) {$5$};
        \node[left, font=\tiny, grisTexto] at (0,9.8) {$Y=10$};

        % Curva de Contrato y_A = x_A y Núcleo [(4,4) a (6,6)]
        \draw[verdeAcento!75, thick, dashed] (0,0) -- (10,10);
        \draw[verdeAcento, line width=2.3pt] (4,4) -- (6,6);

        % Indiferencias por W(8,2)
        \draw[azulPrincipal, thick, domain=1.65:9.2, samples=40] plot (\x, {16/\x});
        \node[azulPrincipal, font=\tiny\bfseries, anchor=west] at (1.75, 9.25) {$U_A=16$};
        \draw[rojoAcento, thick, domain=0.8:8.28, samples=40] plot (\x, {10 - 16/(10-\x)});
        \node[rojoAcento, font=\tiny\bfseries, anchor=south] at (5.6, 6.5) {$U_B=16$};

        % Indiferencias óptimas tangentes en E*(5,5): U_A*=25 y U_B*=25
        \draw[azulPrincipal!85, thick, dashed, domain=2.75:9.2, samples=40] plot (\x, {25/\x});
        \draw[rojoAcento!85, thick, dashed, domain=0.8:7.25, samples=40] plot (\x, {10 - 25/(10-\x)});

        % Recta presupuestaria x + y = 10 (px/py = 1)
        \draw[grisTexto!90!black, thick] (1.0, 9.0) -- (9.0, 1.0);
        \node[grisTexto, font=\tiny, inner sep=0pt, scale=0.88, anchor=east] at (9.90, 0.44) {$p_x/p_y=1$};

        % Puntos clave W=(8,2), E*=(5,5), T=(6,6)
        \fill[naranjaAcento] (8,2) circle (3.5pt) node[anchor=west, font=\tiny\bfseries] {$W(8,2)$};
        \fill[azulOscuro] (5,5) circle (3.6pt);
        \draw[<-, azulOscuro, thin] (5.00, 5.10) -- (4.92, 5.72) node[above=0pt, font=\tiny\bfseries, align=center, inner sep=0.5pt, anchor=south] {$E^*$\\[-2pt]$(5,5)$};
        \fill[verdeAcento!90!black] (6,6) circle (3.0pt) node[anchor=west, font=\tiny\bfseries] {$(6,6)$};
        \node[azulPrincipal, font=\tiny\bfseries, anchor=west] at (3.2, 8.2) {$U_A^*=25$};
      \end{tikzpicture}
    \end{column}

    \begin{column}{0.46\textwidth}
      {\scriptsize\color{grisTexto}Resolución Paso a Paso (Caso Clásico Bloque B)}\\[2pt]
      {\scriptsize
        \textbf{1. Evaluar \(W=(8,2)\):} \(TMS^A = \frac{2}{8} = 0.25 \neq 4 = TMS^B \implies\) intercambio.\\[3pt]
        \textbf{2. Curva de Contrato:} \(\frac{y_A}{x_A} = \frac{10-y_A}{10-x_A} \implies \mathbf{y_A = x_A}\).\\[3pt]
        \textbf{3. Equilibrio (\(E^*\)):} con \(p_y = 1\), \(I_A = 8p_x+2\) y \(x_A + x_B = 10 \implies \mathbf{p_x / p_y = 1}\), \(E^* = (5,5)\), \(\Delta U_A = \Delta U_B = \mathbf{+9}\).\\[3pt]
        \textbf{4. Ajuste (Ley de Walras):} \(p_x z_x + p_y z_y \equiv 0\). Si \(z_x > 0\), entonces \(z_y < 0\) y sube \(p_x / p_y\).\\[4pt]
      }
      \recuadroclave{
        {\scriptsize\textbf{2do Teorema (M04-Q44):} Para lograr \(A=(6,6),\; B=(4,4)\) con \(p_x=p_y=1\), transferir \(T = \mathbf{2}\) de \(B\) a \(A\).}
      }
    \end{column}
  \end{columns}

  \mensajefinal{En el Núcleo de \(W=(8,2)\), ambos superan su utilidad inicial \(U_A=U_B=16\).}
\end{frame}

% ==============================================================================
% SLIDE 4: EFICIENCIA GLOBAL CON PRODUCCIÓN
% ==============================================================================
\begin{frame}{2. Equilibrio General con Producción (3 Condiciones de Pareto)}
  \vspace{0.1cm}
  \begin{columns}[T]
    \begin{column}{0.32\textwidth}
      \tarjeta{azulPrincipal}{1. Eficiencia en el Consumo}{
        \(TMS_{1,2}^A = TMS_{1,2}^B\)\\[4pt]
        Asigna los bienes producidos de forma óptima entre los consumidores (sobre la Curva de Contrato de consumo).
      }
    \end{column}
    \begin{column}{0.32\textwidth}
      \tarjeta{naranjaAcento}{2. Eficiencia en la Producción}{
        \(TMST_{L,K}^X = TMST_{L,K}^Y\)\\[4pt]
        Asigna capital y trabajo de forma óptima entre industrias (sobre la \textbf{Frontera de Posibilidades de Producción - FPP}).
      }
    \end{column}
    \begin{column}{0.32\textwidth}
      \tarjeta{verdeAcento}{3. Eficiencia Global (Mix)}{
        \(TMT_{X,Y} = TMS_{X,Y}^A = TMS_{X,Y}^B\)\\[4pt]
        La tasa a la que la economía transforma \(X\) en \(Y\) (\(\frac{CMg_X}{CMg_Y}\)) iguala la valoración marginal de los consumidores.
      }
    \end{column}
  \end{columns}

  \vspace{0.2cm}
  \recuadroclave{
    \small\textbf{FPP y Costo de Oportunidad:} La pendiente de la FPP es \(-TMT_{X,Y} = -\frac{CMg_X}{CMg_Y}\). En competencia perfecta: \(\frac{p_X}{p_Y} = \frac{CMg_X}{CMg_Y} = TMT_{X,Y}\).
  }
  \mensajefinal{Producir sobre la FPP garantiza eficiencia técnica, pero solo igualar \(TMT = TMS\) garantiza eficiencia económica global.}
\end{frame}

% ==============================================================================
% SLIDE 5: LOS DOS TEOREMAS DEL BIENESTAR
% ==============================================================================
\begin{frame}{3. Los Dos Teoremas Fundamentales del Bienestar}
  \vspace{0.15cm}
  \begin{columns}[T]
    \begin{column}{0.48\textwidth}
      \tarjeta{azulPrincipal}{Primer Teorema del Bienestar}{
        \textbf{Enunciado:} Todo equilibrio competitivo (walrasiano) es un \textbf{Óptimo de Pareto}.\\[4pt]
        • \textbf{Supuesto clave:} Preferencias localmente no saciadas (y mercados completos, sin externalidades ni poder de mercado).\\
        • \textbf{No requiere convexidad.}\\
        • \textit{Significado:} La ``mano invisible'' logra eficiencia, pero \textbf{no garantiza equidad}.
      }
    \end{column}
    \begin{column}{0.48\textwidth}
      \tarjeta{verdeAcento}{Segundo Teorema del Bienestar}{
        \textbf{Enunciado:} Cualquier Óptimo de Pareto deseado puede sostenerse como equilibrio competitivo mediante \textbf{transferencias de suma alzada} (redistribución inicial) y dejando actuar al mercado.\\[4pt]
        • \textbf{Supuesto clave:} \textbf{Preferencias y conjuntos de producción convexos.}\\
        • \textit{Significado:} Separa el problema de la \textbf{eficiencia} del problema de la \textbf{equidad}.
      }
    \end{column}
  \end{columns}

  \vspace{0.2cm}
  \recuadroclave{
    \small\textbf{Lección de Política Regulatoria:} Para ayudar a hogares pobres en agua potable sin distorsionar precios, es mejor dar un \textbf{bono directo (suma alzada)} que distorsionar el precio por debajo del \(CMg\).
  }
  \mensajefinal{Los impuestos distorsionadores rompen la igualdad \(TMS = TMT\) y violan el Segundo Teorema.}
\end{frame}

% ==============================================================================
% SLIDE 6: FUNCIONES DE BIENESTAR SOCIAL Y TEOREMA DE ARROW
% ==============================================================================
\begin{frame}{4. Elección Social, Equidad y Teorema de Imposibilidad de Arrow}
  \vspace{0.1cm}
  \begin{columns}[T]
    \begin{column}{0.32\textwidth}
      \tarjeta{azulPrincipal}{Utilitarista (Bentham)}{
        \(W = \sum_{i=1}^n U_i\)\\[4pt]
        Maximiza la suma total de utilidades. Indiferente a cómo se distribuye entre ricos y pobres si la suma crece.
      }
    \end{column}
    \begin{column}{0.32\textwidth}
      \tarjeta{verdeAcento}{Rawlsiana (Maximin)}{
        \(W = \min\{U_1, U_2, \dots, U_n\}\)\\[4pt]
        Maximiza el bienestar del individuo \textbf{peor situado}. Curvas de indiferencia social en forma de ``L'' (Leontief).
      }
    \end{column}
    \begin{column}{0.32\textwidth}
      \tarjeta{rojoAcento}{Teorema de Arrow}{
        No existe regla de votación colectiva que cumpla simultáneamente:\\
        1. Dominio universal\\
        2. Eficiencia de Pareto\\
        3. Independencia de alternativas irrelevantes (IIA)\\
        4. \textbf{No dictadura}
      }
    \end{column}
  \end{columns}

  \vspace{0.15cm}
  \recuadroclave{
    \small\textbf{Paradoja de Condorcet:} La regla de mayoría simple puede generar preferencias sociales \textbf{intransitivas} (ciclos de votación \(A \succ B \succ C \succ A\)).
  }
  \mensajefinal{El óptimo social se alcanza donde la curva de isobienestar social es tangente a la Frontera de Posibilidades de Utilidad (FPU).}
\end{frame}

% ==============================================================================
% SLIDE 7: PRODUCCIÓN 2x2, TEOREMAS H-O Y CRITERIO KALDOR-HICKS
% ==============================================================================
\begin{frame}{5. Estática Comparativa \(2\times 2\), Existencia y Kaldor-Hicks}
  \vspace{0.01cm}
  \begin{columns}[T]
    \begin{column}{0.32\textwidth}
      \tarjeta{azulPrincipal}{Teoremas \(2\times 2\) (Producción)}{
        • \textbf{Stolper-Samuelson:} Si sube \(P_X\), aumenta la retribución real del factor usado \textbf{intensivamente en \(X\)} y cae la del otro factor.\\
        • \textbf{Rybczynski:} Si crece la dotación de \(L\), aumenta más que proporcionalmente la producción del bien intensivo en \(L\) y \textbf{cae} la del otro.
      }
    \end{column}
    \begin{column}{0.32\textwidth}
      \tarjeta{verdeAcento}{Existencia y Unicidad}{
        • \textbf{Existencia (Arrow-Debreu):} Teorema de Punto Fijo de \textbf{Brouwer} (o \textbf{Kakutani}) si \(z(p)\) es continuo y homogéneo.\\
        • \textbf{Unicidad:} Exige sustitubilidad bruta o WARP agregado.
      }
    \end{column}
    \begin{column}{0.32\textwidth}
      \tarjeta{naranjaAcento}{Criterio de Kaldor-Hicks}{
        Un cambio mejora bienestar si ganadores \textbf{podrían compensar} a perdedores y seguir mejor (mejora potencial de Pareto).\\[2pt]
        • No exige compensación efectiva; base del análisis Costo-Beneficio (\(\Delta W > 0\)).
      }
    \end{column}
  \end{columns}

  \vspace{0.02cm}
  \recuadroclave{
    \small\textbf{FPP Cóncava vs. Lineal:} La FPP es cóncava al origen cuando existen intensidades factoriales distintas entre \(X\) e \(Y\) o rendimientos marginales decrecientes; solo es recta si hay proporciones idénticas y RCE.
  }
  \mensajefinal{Sonnenschein-Mantel-Debreu: el equilibrio general puede ser múltiple e inestable.}
\end{frame}

% ==============================================================================
% SLIDE 8: CONCLUSIONES Y CHEAT-SHEET MÓDULO IV
% ==============================================================================
\begin{frame}{Conclusiones y Formulario Rápido — Módulo IV}
  \vspace{0.05cm}
  \recuadroclave{
    \small\textbf{Regla de Oro:} El mercado competitivo alcanza la eficiencia paretiana (\(TMS = TMST = TMT\)), pero la equidad depende de la distribución inicial o transferencias de suma alzada.
  }

  \vspace{0.15cm}
  \begin{columns}[T]
    \begin{column}{0.32\textwidth}
      \tarjeta{verdeAcento}{Fórmulas Imprescindibles}{
        • Ley de Walras: \(\sum p_i z_i(p) \equiv 0\)\\
        • Consumo: \(TMS^A = TMS^B\)\\
        • Global: \(TMT_{X,Y} = \frac{CMg_X}{CMg_Y} = TMS_{X,Y}\)
      }
    \end{column}
    \begin{column}{0.32\textwidth}
      \tarjeta{azulPrincipal}{Trampas de Examen}{
        • Creer que el 1er Teorema requiere convexidad (solo el 2do la exige).\\
        • Confundir Óptimo de Pareto con equidad o justicia.\\
        • Olvidar que un punto dentro del lente mejora a ambos agentes.
      }
    \end{column}
    \begin{column}{0.32\textwidth}
      \tarjeta{rojoAcento}{Siguientes Pasos}{
        • Practicar Bloque A y B del Módulo 04 en el Simulador.\\
        • Repasar cálculo de curva de contrato con Cobb-Douglas.
      }
    \end{column}
  \end{columns}

  \vspace{0.25cm}
  \begin{center}
    {\color{azulPrincipal}\Large\textbf{¡Éxito!}} \quad {\color{grisTexto}Módulo IV · Equilibrio General y Bienestar · PEU SUNASS 2026}
  \end{center}
\end{frame}

\end{document}
